% Section 6 --- Lower Bounds (~1.5 pages, weight 12%)
% 6.1 Theorem 1 --- Necessary Capability for Deterministic A2 (~0.4 page proof)
% 6.2 Observation 1 --- Capability classification
% 6.3 Theorem 2 --- Calibrated Bond Insufficiency (~0.3 page proof)
% 6.4 Corollary --- Empirical violation rate
%
% REVIEW STATUS: Draft. Awaiting Task 8.2 (distributed-systems colleague review).
% OPEN: Theorem 1 proof Step "Construction" --- confirm adversary model
%       is consistent with §2.5 threat model.
% OPEN: Theorem 2 --- C and alpha_hat to be substituted from Task 7.4 results.

\section{Lower Bounds}
\label{sec:bounds}

\subsection{Theorem 1: Necessary Capability for Deterministic A2}
\label{sec:thm1}

We first fix notation. Let $\mathcal{R} = \{R_1, R_2\}$ be two rollups sharing
a sequencer $\mathcal{S}$. Each rollup $R_i$ has state $\sigma_i$ and state
transition function $\delta_i(\sigma_i, \tau_i) \to \sigma_i'$, where $\tau_i$
is a transaction. A \emph{bundle} $B = (\tau_1, \tau_2)$ pairs one transaction
per rollup. The sequencer issues a \emph{commitment} $\Gamma(B)$ attesting to
the inclusion and ordering of $B$ before execution.

\begin{definition}[A2 --- Execution Atomicity]
\label{def:a2}
A sequencer $\mathcal{S}$ provides \emph{A2 execution atomicity} for bundle $B$
if, for every execution of $B$ under any admissible interleaving of $R_1, R_2$:
\[
  \delta_1(\sigma_1, \tau_1) \text{ succeeds} \iff \delta_2(\sigma_2, \tau_2) \text{ succeeds.}
\]
That is, either both legs apply their state transitions, or neither does.
\end{definition}

\begin{definition}[Lazy Sequencer]
\label{def:lazy}
A sequencer $\mathcal{S}$ is \emph{lazy} if it issues $\Gamma(B)$ without
evaluating $\delta_1$ or $\delta_2$. Formally: the computation of $\Gamma(B)$
does not depend on the output of $\delta_i(\sigma_i, \tau_i)$ for either $i$.
\end{definition}

\begin{definition}[Independent STFs]
\label{def:independent}
State transition functions $\delta_1, \delta_2$ are \emph{independent} if
$\delta_1(\sigma_1, \tau_1)$ does not read or modify $\sigma_2$, and
$\delta_2(\sigma_2, \tau_2)$ does not read or modify $\sigma_1$.
\end{definition}

\begin{definition}[Irreversibility Point]
\label{def:irreversibility}
For rollup $R_i$, the \emph{irreversibility point} of a transaction $\tau_i$
is the first protocol-internal moment at which the effect of $\tau_i$ on
$\sigma_i$ cannot be undone by any party within the protocol. Denote this
moment $\pi_i(\tau_i)$.
\end{definition}

\begin{theorem}[Necessary Capability for Deterministic A2]
\label{thm:a2-necessary}
Let $\mathcal{S}$ be a deterministic sequencer that is lazy
(Definition~\ref{def:lazy}) and sequences bundles over independent STFs
(Definition~\ref{def:independent}). If $\mathcal{S}$ provides A2 execution
atomicity (Definition~\ref{def:a2}), then $\mathcal{S}$ must implement at
least one of the following capabilities:
\begin{enumerate}
  \item[(i)] \textbf{Pre-execution evidence:} Before issuing $\Gamma(B)$,
    $\mathcal{S}$ obtains a success predicate $\phi(B)$ such that
    $\phi(B) = \mathsf{true}$ implies both $\delta_1(\sigma_1, \tau_1)$ and
    $\delta_2(\sigma_2, \tau_2)$ will succeed under any admissible execution.
  \item[(ii)] \textbf{Partial-effect prevention:} $\mathcal{S}$ has authority
    to prevent the irreversible effect of $\tau_i$ on $R_i$ from persisting if
    $\delta_j(\sigma_j, \tau_j)$ fails for $j \neq i$. Formally, for any
    execution where $\delta_1$ succeeds at $\pi_1(\tau_1)$ but $\delta_2$
    subsequently fails, $\mathcal{S}$ can trigger a state revert of $R_1$ to
    $\sigma_1$ before $\pi_1(\tau_1)$ is externally observable.
\end{enumerate}
\end{theorem}

\begin{proof}
By contrapositive. Assume $\mathcal{S}$ satisfies neither (i) nor (ii); we
construct an execution that violates A2.

\medskip
\noindent\textbf{Setup.} Since $\mathcal{S}$ is lazy, it issues $\Gamma(B)$
without evaluating $\delta_1$ or $\delta_2$. Since (i) is not satisfied, there
is no success predicate available to $\mathcal{S}$ at commitment time. Therefore,
$\Gamma(B)$ may be issued even when one or both of $\delta_1(\sigma_1, \tau_1)$,
$\delta_2(\sigma_2, \tau_2)$ would fail.

\medskip
\noindent\textbf{Construction of counterexample execution.} Because $\delta_1,
\delta_2$ are independent (Definition~\ref{def:independent}), their outcomes
depend only on their respective local states $\sigma_1, \sigma_2$. Consider
the following scenario. An adversary $\mathcal{A}$ holds a bundle
$B = (\tau_1, \tau_2)$, submits $B$ to $\mathcal{S}$, and obtains $\Gamma(B)$.
Between $\Gamma(B)$ and the execution of $\tau_2$ on $R_2$, $\mathcal{A}$
submits a conflicting transaction $\tau_2'$ to $R_2$ that exhausts the resource
$\tau_2$ depends on, causing $\delta_2(\sigma_2, \tau_2)$ to fail.

Such a $\tau_2'$ exists because: (a) $R_2$'s state evolves independently of
$R_1$ (by Definition~\ref{def:independent}), so $\tau_2'$ can target $\sigma_2$
without being blocked by any constraint on $R_1$; and (b) $\mathcal{S}$ is lazy
and does not serialize access to $R_2$'s state between $\Gamma(B)$ and execution.

\medskip
\noindent\textbf{Result without capability (ii).} In the constructed execution,
$\delta_1(\sigma_1, \tau_1)$ succeeds (its preconditions are unaffected by
$\tau_2'$ due to independence) and reaches its irreversibility point
$\pi_1(\tau_1)$. Meanwhile, $\delta_2(\sigma_2, \tau_2)$ fails due to
$\tau_2'$. Since $\mathcal{S}$ does not have capability (ii), it cannot revert
the effect of $\tau_1$ on $R_1$ after $\pi_1(\tau_1)$.

\medskip
\noindent\textbf{A2 violation.} Definition~\ref{def:a2} requires
$\delta_1$ succeeds $\iff$ $\delta_2$ succeeds. In the constructed execution,
$\delta_1$ succeeds and $\delta_2$ fails. This directly violates A2. \qed
\end{proof}

\noindent\textbf{Scope remark.} Theorem~\ref{thm:a2-necessary} applies to
\emph{deterministic} protocols. Probabilistic and cryptoeconomic guarantees ---
including bond-slashing and restaking --- do not provide either (i) or (ii).
A bond guarantees that if A2 fails, the violating party loses stake; it does not
prevent the A2 failure itself. This is the A4 guarantee, treated separately in
Theorem~\ref{thm:bond-insufficient} and \S\ref{sec:obs1}.

%-----------------------------------------------------------------------

\subsection{Observation 1: Capability Classification of Deployed Mechanisms}
\label{sec:obs1}

Theorem~\ref{thm:a2-necessary} partitions deployed mechanisms into three classes.
Table~\ref{tab:capability-class} classifies the mechanisms surveyed in
\S\ref{sec:designspace}; boundary cases are analyzed in Appendix~\ref{app:boundary}.

\begin{table}[h]
\centering
\small
\caption{Capability classification. Mechanisms providing only A4 guarantees
  cannot close the A2 gap regardless of bond size. See Appendix~\ref{app:boundary}
  for boundary case rationale.}
\label{tab:capability-class}
\begin{tabular}{@{}lll@{}}
\toprule
Mechanism & Capability & Basis \\
\midrule
ZK proof gating           & (i)       & Commitment requires execution proof \\
Eager simulation          & (i)       & Sequencer runs STF before committing \\
Espresso exec-attest      & (i)       & HotShot execution attestation \\
HTLC                      & (ii)      & Hashlock prevents partial finality \\
Across Protocol escrow    & (ii)      & Funds locked until both legs confirm \\
Astria bundle escrow      & (ii)      & App-level escrow over sequencer bundle \\
Shutter encryption        & (i)+(ii)  & Threshold decryption prevents selective exercise \\
Bond / slashing           & A4 only   & Compensates; does not prevent A2 failure \\
EigenLayer restaking      & A4 only   & Extends slashing domain; same execution model \\
MEV-Boost relay attest    & A4 only   & Economic attestation; no execution proof \\
Bonded preconf (mev-commit, Bolt) & A4 only & Bond on inclusion; no STF execution \\
\bottomrule
\end{tabular}
\end{table}

%-----------------------------------------------------------------------

\subsection{Theorem 2: Calibrated Bond Insufficiency}
\label{sec:thm2}

We show that finite bonds cannot eliminate A4 violations under a power-law
option-value tail, even when the bond is set optimally.

\begin{definition}[A4 violation indicator]
\label{def:a4-viol}
For a commitment $c$ in window $t$, define the A4 violation indicator:
\[
  \mathbf{1}_{A4}(c_t) \;=\; \mathbf{1}\!\left[V(c_t) > B\right]
\]
where $V(c_t)$ is the option value of the commitment and $B$ is the sequencer's
bond. A violation occurs when exercising the option is more profitable than
forfeiting the bond.
\end{definition}

\begin{theorem}[Bond Insufficiency under Power-Law Tail]
\label{thm:bond-insufficient}
Let $V$ be the option value of a sequencing commitment with power-law tail:
$\Pr[V > v] \sim C v^{-\alpha}$ for $v \to \infty$, tail index $\hat{\alpha} > 0$,
$C > 0$. Let $B_{\max}$ be the maximum feasible bond, $w$ the number of
measurement windows, and $\lambda$ the per-window commitment arrival rate.
Then:
\[
  \mathbb{E}\!\left[\sum_{t=1}^{w} \mathbf{1}_{A4}(c_t)\right]
  \;\geq\;
  \rho_{\min}(\hat{\alpha},\, B_{\max},\, w,\, \lambda)
  \;\triangleq\;
  w \cdot \lambda \cdot C \cdot B_{\max}^{-\hat{\alpha}} \cdot (1-\varepsilon)
\]
for any $\varepsilon > 0$. The bound $\rho_{\min}$ is strictly positive for any
finite $B_{\max}$, any $\hat{\alpha} > 0$, $w > 0$, $\lambda > 0$.

\noindent\textbf{Remark.} The requirement $\hat{\alpha} > 0$ is weaker than the
common assumption $\hat{\alpha} > 1$ (finite mean). When $0 < \hat{\alpha} \leq 1$,
$V$ has infinite mean, making the bound even stronger: tail events occur with
probability decaying as $B_{\max}^{-\hat{\alpha}}$, which diverges faster as
$\hat{\alpha} \to 0$. Chains with $\hat{\alpha} < 1$ (e.g., Arbitrum, where
congestion creates super-heavy-tailed fee spikes) exhibit even higher expected
violation rates than chains with $\hat{\alpha} > 1$.
\end{theorem}

\begin{proof}
For each window $t$, by the power-law tail assumption:
\[
  \Pr\!\left[V_t > B_{\max}\right] \;\geq\; C \cdot B_{\max}^{-\hat{\alpha}} \cdot (1-\varepsilon)
\]
for sufficiently large $B_{\max}$ (and for all $B_{\max}$ if the inequality
holds exactly). Since $V_t > B_{\max}$ implies $\mathbf{1}_{A4}(c_t) = 1$:
\[
  \Pr[\mathbf{1}_{A4}(c_t) = 1] \;\geq\; C \cdot B_{\max}^{-\hat{\alpha}} \cdot (1-\varepsilon).
\]
There are $w \lambda$ commitment windows in expectation. By linearity of
expectation (windows are independent by the system model in \S\ref{sec:model}):
\[
  \mathbb{E}\!\left[\sum_t \mathbf{1}_{A4}(c_t)\right]
  \;\geq\; w \cdot \lambda \cdot C \cdot B_{\max}^{-\hat{\alpha}} \cdot (1-\varepsilon). \qed
\]
\end{proof}

The derivation of $\rho_{\min}$ in closed form, with empirical $\hat{C}$ and
$\hat{\alpha}$ from the 12-month replay, is in Appendix~\ref{app:theorem2}.

%-----------------------------------------------------------------------

\subsection{Corollary: Empirical Violation Rate}
\label{sec:corollary}

We calibrate Theorem~\ref{thm:bond-insufficient} with empirical parameters
from the 12-month counterfactual replay (\S\ref{sec:meas-replay}).

\paragraph{Chain-specific calibration.}
A key finding of the measurement study is that the \emph{relevant} bond size
$B_{\max}$ is chain-dependent: the option value distribution $V$ varies
substantially across L2 chains, so a bond that suffices on one chain may be
inadequate on another.

Table~\ref{tab:calibrated-rhomin} reports the empirical A4 violation rate
$\hat{\rho}(B_{\max})$ and the Theorem~\ref{thm:bond-insufficient} lower bound
$\rho_{\min}(\hat{\alpha}, B_{\max})$ at chain-calibrated thresholds.
The calibrated $B_{\max}$ for each chain is the 90th percentile of $V$
on that chain --- the bond level at which the top 10\% of blocks would
trigger violations.

\begin{table}[h]
\centering
\small
\caption{Empirical A4 violation rate vs.\ Theorem~2 bound at
  chain-calibrated $B_{\max}$ (90th percentile of $V$).
  $\hat{\rho}$: fraction of replay windows with $V > B_{\max}$.
  $\rho_{\min}$: Theorem~2 lower bound.
  Bold: mev-commit reference scenario ($B_{\max} = 0.05$\,ETH).}
\label{tab:calibrated-rhomin}
\begin{tabular}{@{}llrrr@{}}
\toprule
Chain & $B_{\max}$ & $\hat{\rho}_{\text{emp}}$ & $\rho_{\min}$ (T2) & $\hat{\rho}/\rho_{\min}$ \\
\midrule
Optimism  & 0.1\,mETH    & 29.0\% & 27.8\% & 1.04$\times$ \\
Optimism  & 1\,mETH      &  2.0\% &  1.2\% & 1.63$\times$ \\
Base      & 10\,mETH     & 18.0\% & 17.5\% & 1.03$\times$ \\
\textbf{Base} & \textbf{50\,mETH (0.05 ETH)} & \textbf{2.0\%} & \textbf{2.6\%} & \textbf{0.77$\times$} \\
Arbitrum  & 10\,mETH     & 10.7\% & 10.7\% & 1.00$\times$ \\
Arbitrum  & 50\,mETH     &  6.4\% &  5.0\% & 1.28$\times$ \\
\bottomrule
\end{tabular}
\end{table}

\paragraph{Key observations.}

\textbf{(a) T2 bound is tight and consistent.}
At calibrated thresholds, the Theorem~2 bound $\rho_{\min}$ tracks the
empirical rate within a factor of $0.77$--$1.63\times$ across all chains.
The bound is an \emph{asymptotic} lower bound (formally valid as
$B_{\max} \to \infty$); at finite $B_{\max}$, the empirical rate can be
slightly below $\rho_{\min}$ due to the approximation error in the
power-law tail estimate~\cite{Clauset2009}. The ratio of 0.77$\times$ at
Base~$B_{\max} = 0.05$\,ETH is within the expected range of finite-sample
estimation noise (bootstrap CI on $\hat{\alpha}$ spans $[1.00, 1.48]$).
Over the full sweep, the bound is neither vacuous nor systematically
biased.

\textbf{(b) B\textsubscript{max} calibration is chain-dependent.}
At the mev-commit reference scenario ($B_{\max} = 0.05$\,ETH):
\begin{itemize}
  \item \textbf{Optimism}: empirical rate $= 0\%$ (not because the bound is
    vacuous, but because 50\,mETH far exceeds the typical Optimism block fee
    of 138\,$\mu$ETH; the relevant calibrated threshold for Optimism is
    0.1--1\,mETH, where $\hat{\rho} = 2$--$29\%$).
  \item \textbf{Base}: empirical rate $= 2.0\%$, T2 bound $= 2.6\%$.
    Base's higher-value blocks (mean 9.8\,mETH/block) place 0.05\,ETH within
    the observable tail. This is the most directly comparable scenario to
    the mev-commit worked example (\S\ref{sec:peo-example}).
  \item \textbf{Arbitrum}: empirical rate $= 6.4\%$, T2 bound $= 5.0\%$.
    Arbitrum's super-heavy-tailed burst-congestion events make even large
    bonds regularly exceeded.
\end{itemize}

\textbf{(c) All-chains confirmation.}
Across all chains and all calibrated thresholds, $\hat{\rho} > 0$ and
$\rho_{\min} > 0$. No finite bond eliminates A4 violations on any of the
measured chains.

\paragraph{Primary calibration statement.}
Using the Base chain at $B_{\max} = 0.05$\,ETH (direct mev-commit comparison):
\[
  \hat{\rho}_{\text{emp}}(\text{Base},\; 0.05\,\text{ETH}) = 2.0\%
  \quad \geq \quad
  \rho_{\min}(1.19,\; 0.05,\; w,\; \lambda) = 2.6\%
  \;\text{(T2 bound).}
\]
Over a 12-month window with $\lambda = 1$ commitment per block and
$w = 15{,}768{,}000$ Base blocks/year, the expected number of A4
violations at the mev-commit reference bond is bounded below by
$\rho_{\min} \times w = 0.026 \times 15.8\text{M} \approx 410{,}000$
expected violations per year --- a structural, not incidental, consequence
of the power-law tail.

\noindent\textbf{Interpretation.}
The calibrated violation rates confirm that no finite bond eliminates A4
violations. The appropriate bond to suppress violations to a given rate is
chain-dependent and must be calibrated to the actual option value distribution:
a bond adequate for Optimism (where 50\,mETH $\gg$ typical block fee) would
provide substantially weaker protection on Base or Arbitrum. In all cases,
the fundamental limitation is structural: bonds are A4 mechanisms, not A2
mechanisms. The only path to A2 is capability~(i) or~(ii) per
Theorem~\ref{thm:a2-necessary}.
